\subsection{Movimento orbital relativo}

O movimento relativo $s$ é definido a seguir e as acelerações relativas se tornam diretamente as derivadas no espaço inercial \cite{fehse2003}.

\begin{equation}
\vec r_t + \vec s = \vec r_c
\end{equation}

\begin{equation}
\label{eq:vetorS}
\vect s = \vect r_{c} - \vect r_{t}
\end{equation}

\begin{equation}
\label{eq:vetorS_ddot}
\ddot{\vect s} = \ddot{\vect r}_{c} - \ddot{\vect r}_{t}
\end{equation}

Inserindo as equações \eqref{eq:movimentoChaser} e \eqref{eq:movimentoTarget} do movimento relativo do \emph{chaser} e do \emph{target} na equação \eqref{eq:vetorS_ddot}, temos:

\begin{equation}
\label{eq:vetorS_ddot_substituicao}
\ddot{\vec s} = \left(-\mu \frac{\vect r_c}{\vect r_{c}^3} + \frac{\vect F}{m_c}\right) - \left(-\mu \frac{\vect r_t}{\vect r_t^3}\right)
\end{equation}

Para o próximo passo, serão considerados:

\begin{equation}
    \label{eq:fg_rc}
\vec f_g(\vec r_c) = -\mu \frac{\vect r_c}{\vect r_{c}^3}
\end{equation}

\begin{equation}
    \label{eq:fg_rt}
\vec f_g(\vec r_t) = -\mu \frac{\vect r_t}{\vect r_{t}^3}
\end{equation}

Portanto, a equação \eqref{eq:vetorS_ddot_substituicao} se torna:

\begin{equation}
\label{eq:fg_rc_fg_rt}
\ddot{\vect s} = \vec f_g(\vect r_c) - \vec f_g(\vect r_t) + \frac{\vect F}{m_c}
\end{equation}

Levando em consideração apenas a força de controle e linearizando a força gravitacional no \emph{chaser} em torno da posição do \emph{target}, se faz necessário realizar uma expansão de Taylor apenas para a primeira ordem, e é obtido \cite{arantes2011}:

\begin{equation}
\label{eq:linearizandoJacobiana}
\vec f(\vec r_c) - \vec f(\vec r_t) = \nabla \vec f(\vec r_t)\vec p
\end{equation}

Os termos da diagonal principal da matriz Jacobiana \emph{J} são dados por \cite{fehse2003}:

\begin{equation}
\label{eq:jacobianaDiagonal}
J_{ii} = \frac{\partial f_i}{\partial r_i} = -\frac{\mu}{r^3} \left(1 - 3\frac{r_i^2}{r^2}\right) \quad (i = j)
\end{equation}

E os termos fora da diagonal principal são dados por \cite{fehse2003}:
\begin{equation}
\label{eq:jacobianaForaDiagonal}
J_{ij} = \frac{\partial f_i}{\partial r_i} = -3 \frac{\mu}{r^3} \frac{r_i r_j}{r^2} \quad (i \neq j)
\end{equation}

Portanto, a matriz Jacobiana \emph{J} será:

\begin{equation}
    \label{eq:jacobianaGenerica}
\vec J = 
\begin{bmatrix}
1 - 3\frac{r_x^2}{r^2} &  -3\frac{r_x r_y}{r^2} &  -3\frac{r_x r_z}{r^2} \\
-3\frac{r_y r_x}{r^2} & 1 - 3\frac{r_y^2}{r^2} & -3\frac{r_y r_z}{r^2} \\
-3\frac{r_z r_x}{r^2} & -3\frac{r_z r_y}{r^2} & 1 - 3\frac{r_z^2}{r^2}
\end{bmatrix}
\end{equation}

Reescrevendo a equação \eqref{eq:linearizandoJacobiana} com os termos da matriz Jacobiana, temos:

\begin{equation}
\label{eq:linearizandoJacobiana2}
\vec f_g(\vec r_c) - \vec f_g(\vec r_t) = -\frac{\mu}{r_t^3} Js
\end{equation}

Portanto, a equação \eqref{eq:fg_rc_fg_rt} se torna:

\begin{equation}
\label{eq:fg_rc_fg_rt_linearizada}  
\ddot{\vect s} = \frac{\mu}{r_t^3} Js + \frac{\vect F}{m_c}
\end{equation}

Para realizar a transformação do \emph{quadro inercial} para o \emph{quadro do target}, é necessário aplicar o \emph{Teorema do Transporte}, que é uma forma de transformar as equações do movimento de um sistema de referência para outro que está em movimento relativo \cite{wie2008space, symon1979}. Para a mesma posição relativa entre o \emph{chaser} e o \emph{target}, será considerado que $\mathbf s = \mathbf s^{*}$.

\begin{equation}
\label{eq:TransporteSymon1979}
\ddot{\mathbf s} = 
\ddot{\mathbf s}^{*}
+ 2\boldsymbol\omega\times\dot{\mathbf s}^{*}
+ \boldsymbol\omega\times\big(\boldsymbol\omega\times\mathbf s^{*}\big)
+ \dot{\boldsymbol\omega}\times\mathbf s^{*}
\end{equation}

A equação \eqref{eq:fg_rc_fg_rt_linearizada} que está no quadro inercial ao ser transformada para o quadro LVLH (Local Vertical Local Horizontal) se torna \cite{fehse2003}:

\begin{equation}
\label{eq:relativo_LVLHgeral}
\ddot{\mathbf s}^{*}
+ 2\,\boldsymbol\omega\times\dot{\mathbf s}^{*}
+ \boldsymbol\omega\times\big(\boldsymbol\omega\times\mathbf s^{*}\big)
+ \dot{\boldsymbol\omega}\times\mathbf s^{*}
- \frac{\mu}{r_t^3}\,\mathbf J\,\mathbf s^{*}
= \frac{\mathbf F}{m_c}
\end{equation}

%%%%% Aqui acaba mov. orbital relativo

% Quando expresso no quadro do \emph{target}, é obtido seu vetor posição:

% \begin{equation}
% \label{eq:posicao_LVLH_target}
% \vec r_t = \begin{bmatrix} 0 \\ 0 \\ -r \end{bmatrix}
% \end{equation}

% E também o vetor de velocidade angular:

% \begin{equation}
% \label{eq:velocidadeAngular_LVLH_target}
% \boldsymbol\omega = \begin{bmatrix} 0 \\ -\omega \\ 0
% \end{bmatrix}
% \end{equation}

% Em termos da equação \eqref{eq:relativo_LVLHgeral}, é obtido:

% \begin{equation}
% \label{eq:velAngular_cross_s}
% \boldsymbol\omega \times \boldsymbol s^* =
% \begin{bmatrix}
%     -\omega z \\
%     0         \\
%     \omega x
% \end{bmatrix} 
% \end{equation}

% \begin{equation}
% \label{eq:velAngular_cross_cross_s}
% \boldsymbol\omega \times (\boldsymbol\omega \times \boldsymbol s^*) =
% \begin{bmatrix}
%     -\omega^2 x \\
%     0           \\
%     -\omega^2 z 
% \end{bmatrix}
% \end{equation}

% \begin{equation}
% \label{eq:omega_cross_dsdt}
% \boldsymbol\omega \times \frac{d^* s^*}{dt} =
% \begin{bmatrix}
%     -\omega \dot z \\
%     0                 \\
%     \omega \dot x
% \end{bmatrix}
% \end{equation}

% \begin{equation}
% \label{eq:dOmega_cross_s}
% \frac{d\boldsymbol\omega}{dt} \times \boldsymbol s^* =
% \begin{bmatrix}
%     -\dot\omega z \\
%     0              \\
%     \dot\omega x
% \end{bmatrix}
% \end{equation}

% \begin{equation}
% \label{eq:jacobiana_LVLH}
% \mathbf Js^* =
% \begin{bmatrix}
%     1 & 0 & 0 \\
%     0 & 1 & 0  \\
%     0 & 0 & -2
% \end{bmatrix} \boldsymbol s^* =
% \begin{bmatrix}
%     x \\
%     y \\
%     -2z
% \end{bmatrix}
% \end{equation}

% No caso da órbita circular ($e=0$), a velocidade angular
% \footnote{Os autores \cite{fehse2003,sidi1997spacecraft} definem a \emph{velocidade angular} (ou \emph{angular rate} ou até mesmo \emph{angular velocity}) como $\omega$, outros autores, como \cite{bate1971}, a definem como $\eta$} 
% $\omega$ é constante e coincide com o movimento médio (\emph{mean motion}, $\eta$), definido por \cite{fehse2003}:

% \begin{equation}
% \label{eq:omega_constante}
% \boldsymbol\omega^2 = \frac{\mu}{r_t^3}
% \end{equation}

% Ou podendo ser reescrita como:

% \begin{equation}
% \label{eq:omega_constante2}
% \eta^2 = \frac{\mu}{r_t^3}
% \end{equation}

% Como $r_t=a$, a velocidade angular $\omega^{2}= \eta^2 = \mu/r_t^{3}$, portanto, neste caso em particular de órbita circular, $\omega=n$.

% Então, $\dot{\omega}=0$, e ao inserir a equação \eqref{eq:omega_constante} nos termos da equação \eqref{eq:relativo_LVLHgeral}, é possível obter as equações linearizadas gerais do movimento relativo, conhecidas como Equações de Hill \cite{hill1878,fehse2003, wie2008space}.

% \begin{equation}
%     \label{eq:equacoes_HillX}
% \ddot{x} - 2\omega \dot{z} = \frac{F_x}{m_c}
% \end{equation}

% \begin{equation}
%     \label{eq:equacoes_HillY}
% \ddot{y} + \omega^2 y = \frac{F_y}{m_c}
% \end{equation}

% \begin{equation}
%     \label{eq:equacoes_HillZ}
% \ddot{z} + 2\omega \dot{x} - 3 \omega^2 z = \frac{F_z}{m_c}
% \end{equation}